\chapter{Executor}
\label{chap:executor}

\section{Role}

This chapter covers the host executor in
\filepath{crates/sitas-core/src/executor}. It is a small, single-threaded async
kernel whose wakers, task queues, timers, cancellation, readiness, and
completion dispatch remain visible. The portable \texttt{no\_std}
\texttt{ShardExecutor} is a separate, narrower implementation of the shared
runtime contract.

\section{Why a custom executor?}

The custom executor makes placement part of task semantics. Each task belongs
to one shard, scheduling group, ready queue, and reactor.

The implementation exposes how a waker requeues a task, a dropped timeout
removes timer state, a TCP future records fd interest, and a completion future
keeps buffers alive.

\section{Responsibilities}

\begin{itemize}
    \item Poll ready tasks with a budget so self-waking tasks cannot starve
          timers and I/O forever.
    \item Implement custom task wakers that enqueue tasks onto the local ready
          queue.
    \item Catch task panics at executor boundaries.
    \item Preserve panic payloads where join handles and \texttt{block\_on}
          make them observable.
    \item Provide \texttt{sleep}, \texttt{timeout}, \texttt{race}, and
          \texttt{yield\_now}.
    \item Provide typed join handles and abort support.
    \item Provide cooperative stop tokens, \texttt{Notify}, and
          \texttt{TaskScope}.
    \item Drive readiness futures for read and write interests.
    \item On Linux, drive executor-owned \texttt{io\_uring} read/write-at
          futures.
\end{itemize}

\section{Task lifecycle}

A spawned future is stored inside a task. The task has a stable executor-local
\texttt{TaskId}, an optional name, a scheduling group id, and internal state
describing whether it is queued, polling, waiting, completed, or cancelled.

The lifecycle is explicit:

\begin{center}
\begin{tikzpicture}[
    state/.style={draw=sitasMuted, fill=sitasCodeBackground,
        minimum width=50mm, minimum height=8mm, text width=36mm,
        align=center, font=\footnotesize\ttfamily, rounded corners=1.5pt},
    branch/.style={draw=sitasMuted, fill=white,
        minimum width=50mm, minimum height=8mm, text width=30mm,
        align=center, font=\footnotesize\ttfamily, rounded corners=1.5pt},
    arr/.style={->, >=stealth, thick, color=sitasAccent},
    lbl/.style={font=\tiny\itshape, color=sitasMuted}
]
\node[state] (spawn) at (0, 0) {spawn};
\node[state] (create) at (0, -1.15) {create Task with pinned Future};
\node[state] (queue) at (0, -2.3) {queue Task in scheduling group};
\node[state] (poll) at (0, -3.45) {poll Future};

\node[branch] (ready) at (-3.8, -4.85) {Ready};
\node[state] (completed) at (-3.8, -6.0) {completed};

\node[branch] (pending) at (3.8, -4.85) {Pending};
\node[state] (waiting) at (3.8, -6.0) {waiting until waker fires};
\node[state] (wake) at (3.8, -7.15) {wake};

\draw[arr] (spawn) -- (create);
\draw[arr] (create) -- (queue);
\draw[arr] (queue) -- (poll);
\draw[arr] (poll) -- node[pos=0.50,  left, lbl] {returns} (ready);
\draw[arr] (ready) -- (completed);
\draw[arr] (poll) -- node[pos=0.50,  right, lbl] {returns} (pending);
\draw[arr] (pending) -- (waiting);
\draw[arr] (waiting) -- (wake);
\draw[arr] (wake.east) -- ++(+5mm, 0) |- node[pos=0.25, right, lbl] {requeue} (queue.east);
\end{tikzpicture}
\end{center}

The task records whether it is queued, so repeated wakeups coalesce into one
ready-queue entry.

Cancellation is normal future dropping. Aborting a join handle or dropping an
executor drops the owned future. Timer and readiness futures must clean up
their registrations in \texttt{Drop}; otherwise the reactor keeps stale
interests for nonexistent tasks.

\begin{notabene}[Why cancellation is treated as ordinary]
Rust futures are cancelled by being dropped. The dropped future, task, join
handle, or dispatcher must therefore own the corresponding cleanup.
\end{notabene}

\section{Wakers and the ready queue}

The custom waker does one job: make the task runnable again on the same
executor. It does not move the task to another shard, poll it directly, or
borrow application state. It marks the task queued and hands it to the
scheduler.

The executor polls only a fixed batch of ready work per loop. The
\texttt{READY\_POLL\_BUDGET} constant defines the limit. A self-waking task can
remain ready but cannot prevent checks for timers, readiness, and completions.

The ready queue is executor-local. Cross-thread wakeups may notify an executor
that work became ready, but the task is still polled by its owning executor.
This keeps local task state, scheduling group accounting, and reactor
registrations aligned.

\section{Scheduling groups}

Scheduling groups provide cooperative resource classes. Each group owns a
ready queue and has a relative share count. Ordinary spawn calls use the
default group. Explicit group APIs classify work such as foreground request
handling and background maintenance.

A shard can divide polling time among work classes without moving application
state or introducing a global queue.

The selection rule is simple. The scheduler chooses a non-empty group with the lowest weighted virtual runtime, pops one task from that
group's FIFO queue, polls it, and charges the group for the observed wall-clock
poll duration. The charge is scaled by the group's shares.

\begin{lstlisting}[language={},caption={Weighted group selection}]
for each poll:
    group = non_empty_group_with_lowest_virtual_runtime()
    task = group.pop_front()
    duration = poll(task)
    group.virtual_runtime += duration * default_shares / group.shares
\end{lstlisting}

This is not preemption. A long poll blocks the executor thread until the
future returns \texttt{Poll::Pending} or \texttt{Poll::Ready}. The mechanism is
cooperative: futures must yield, await I/O, await timers, or otherwise return
control to the executor.

Scheduling group handles are executor-local. A group created by one executor is
rejected by another. The default group handle is special because it names the
built-in default queue in any executor. The same idea appears in the sharded
runtime as a \texttt{ShardedSchedulingGroup}, which is one executor-local group
per shard plus a runtime identity check.

\subsection{Scheduling group API shape}

\texttt{Spawner::create\_scheduling\_group} creates an executor-local group
with a name and a positive share count. The normal spawn methods remain the
shorthand for the default group. The group-aware variants classify work before
it enters the ready queues:

\begin{lstlisting}[language={},caption={Local group-aware spawn methods}]
spawn_in_group
spawn_named_in_group
spawn_with_handle_in_group
spawn_with_handle_named_in_group
\end{lstlisting}

\begin{lstlisting}[language={},caption={Classifying local tasks}]
let foreground = spawner.create_scheduling_group("foreground", 200)?;
let background = spawner.create_scheduling_group("background", 25)?;

spawner.spawn_named_in_group(&foreground, "request", request_future)?;
spawner.spawn_named_in_group(&background, "maintenance", maintenance_future)?;
\end{lstlisting}

\texttt{TaskScope} has matching group-aware spawn methods, so lifetime ownership
and scheduling classification can be combined. In the sharded runtime, the
matching creation helper builds one local group per shard:

\begin{lstlisting}[language={},caption={Sharded group creation}]
create_scheduling_group_on_all
\end{lstlisting}

The returned \texttt{ShardedSchedulingGroup} can be used with sharded spawn and
submitter methods that accept an explicit group.

\begin{notabene}[Shares are cooperative]
Shares affect which ready group is selected after tasks return control to the
executor. They do not interrupt a long poll. CPU-bound futures still need to
call \texttt{yield\_now}, await timers or I/O, or split work into bounded
chunks.
\end{notabene}

\section{Structured child tasks}

\texttt{TaskScope} groups child tasks under one owner. Dropping the scope
requests cooperative stop and aborts still-owned children. Calling
\texttt{shutdown} asks children to stop and waits for them. Calling
\texttt{shutdown\_timeout} bounds that wait and aborts uncooperative children.

Scopes can spawn children into scheduling groups, combining lifetime ownership
with resource classification. Scheduling remains the spawner's responsibility.

\section{Wait ordering}

Each loop iteration polls ready work, wakes expired timers, dispatches locally
available \texttt{io\_uring} completions, and then decides how to wait.

On Linux, pending \texttt{io\_uring} submissions are first submitted to the
kernel without blocking. The executor then includes the ring descriptor in the
same reactor wait as read interests, write interests, reactor wakes, and the
next timer deadline. When the ring descriptor becomes readable, the dispatcher
harvests completions and wakes the registered task wakers.

This gives the executor one idle wait shape instead of a separate
\texttt{io\_uring}-first wait path. \texttt{io\_uring} has not become readiness
I/O; completion ownership and buffer lifetime are still handled by the
dispatcher. The reactor only tells the executor that completion queue progress
is available.

\section{Observability}

Executor snapshots are owned values. They include task counts, named task
state, polling counters, timer counts, readiness counts, and Linux
\texttt{io\_uring} dispatcher counters when the dispatcher is installed.
Scheduling group snapshots include ready queue length, total charged poll
count, total charged poll time, weighted virtual runtime, average poll time,
and helper methods for computing each group's share of executor poll time.

Task snapshots include the task name, scheduling group id and name, lifecycle
state, wait interest, poll count, accumulated poll time, and key timestamps.
Helper methods derive task age and current state duration from a caller-supplied
\texttt{Instant}. This keeps the snapshot itself factual while making progress
views easy to write.

\begin{lstlisting}[language={},caption={Reading a task snapshot}]
task name:              useful for humans
status:                 queued, polling, waiting, completed, cancelled
waiting_for:            timer, readable fd, writable fd, unknown, none
poll_count:             how many times the future was polled
total_poll_time:        how much executor time the task consumed
state_duration_at(now): how long it has been in the current state
\end{lstlisting}
